## TEMP holding note — RQ1: does out-of-training-range extrapolation explain XGBoost's edge? (2026-08-19)

**Which claim this checks**: a further candidate reason for why XGBoost beats the DNN on
4survey. Trees split feature space into boxes; each leaf's prediction is bounded by what training
rows landed there. A DNN predicts with a smooth, continuous function that can extrapolate further
from the training data. Under `spatial_block` testing, held-out compartments can have
terrain/environment feature combinations the model never trained on -- a genuine
out-of-training-range generalisation problem, not just interpolation. If this explains part of
XGBoost's edge, the DNN's error should be disproportionately worse on test rows with
out-of-training-range feature values; XGBoost's should not show the same pattern.

**Status: complete (2026-08-19).**

### Method

**Code**: `models/baselines/rq1_extrapolation_check.py` (new file, this session). Run as:
```
PYTHONPATH=. .venv/bin/python -m models.baselines.rq1_extrapolation_check
```

**Pseudocode**:
```
load 4survey, Set3, single spatial_block split (same loader/config as the other RQ1 checks this
    session -- full compartment set, no subsampling)
fit DNN (standard architecture/hyperparameters) and XGBoost (winning 4survey config) on the same
    train/test split, get per-row predictions for both

for each of 14 continuous columns (Age, CanopyCover, time_since_thinning, plus all 11 Set3
    terrain/wind columns) -- excludes binary flags and the categorical thinning_status, which
    have no continuous "out of range" notion:
    compute each test row's distance outside the TRAINING min/max, in units of that column's own
        training SD (so columns on different scales contribute comparably)
sum across columns -> one extrapolation score per test row; also count how many columns are
    out-of-range per row

compare: mean |error| for in-range vs. out-of-range test rows, per model
compare: Spearman correlation between the extrapolation score and |error|, per model
```

### Results

Single-split sanity check: DNN $R^2=0.6166$, XGBoost $R^2=0.6402$ -- consistent with the other
single-split checks this session. 1,936 of 46,956 test rows (4.1\%) have at least one
out-of-training-range feature.

| | In-range (n=45,020) | Out-of-range (n=1,936) | Ratio (out/in) |
|---|---:|---:|---:|
| DNN MAE | 3.5865 | 4.4572 | 1.243 |
| XGBoost MAE | 3.4144 | 4.4782 | 1.312 |

Spearman correlation, extrapolation score vs.\ absolute error: DNN $\rho=0.059$ ($p<10^{-36}$);
XGBoost $\rho=0.074$ ($p<10^{-57}$).

### Headline finding

**Not supported -- if anything, mildly the opposite of what the hypothesis predicts.** Both
models degrade on out-of-range rows (expected -- these are genuinely harder rows for anyone), but
XGBoost's relative degradation is slightly LARGER than the DNN's (ratio 1.312 vs.\ 1.243), and its
correlation with the extrapolation score is slightly STRONGER (0.074 vs.\ 0.059), not weaker. The
"tree leaf predictions stay bounded, so XGBoost handles unseen terrain combinations better" story
does not hold up here -- on this evidence, XGBoost is not more robust to out-of-range test rows
than the DNN is.

**What this does and doesn't establish**: both correlations are weak in absolute terms
($\rho<0.08$), and the out-of-range group is small (4.1\% of test rows), so this is not a
high-powered test of extreme extrapolation. It does rule out this specific mechanism as an
explanation for the *typical* XGBoost-DNN gap, since only a small minority of test rows are
affected at all, and even on those rows XGBoost shows no advantage. A larger-scale distribution
shift (e.g. testing on a genuinely different forest, not just held-out compartments within
Aberfoyle) was not tested and could behave differently.
